\chapter{Mac Lane valuations}\label{ch:mac-lane-valuations}

In his seminal work \cite{MacLane}, S.~Mac Lane considered and essentially solved the following problem.
Let $v_K$ be a nontrivial discrete valuation on a field $K$.
What are the extensions of $v_K$ to an (arbitrary) pseudovaluation $v$ on the polynomial ring $K[x]$?
How can we describe such a pseudovaluation explicitly?
These results have been applied to the problem of constructing models of curves in J.~Rüth's thesis~\cite{Rueth}.
This chapter gives a brief overview of pseudovaluations and key polynomials, augmentations and inductive valuations, minimal augmentation chains, and their relation to discs and discoids.
These constructions provide an explicit counterpart to the geometric valuations used in Part~\ref{part:models-and-valuations}.

Fix an algebraic closure $\overline K$ of $K$ and an extension of $v_K$ to $\overline K$, which we again denote by $v_K$.


\section{Pseudovaluations and key polynomials}\label{sec:pseudovaluations-and-key-polynomials}
We first recall pseudovaluations on $K[x]$ and then single out the Mac Lane pseudovaluations and their key polynomials used in the subsequent constructions.

\begin{definition} \label{def:pseudo-valuation}
    We call $v: K[x] \to \hat{\RR}$ a pseudovaluation if $v(0) = \infty$, $v(1) = 0$, and
    \begin{enumerate}[label=\textup{(\alph*)}]
        \item $v(fg) = v(f) + v(g)$,
        \item $v(f + g) \geq \min\lbrace v(f), v(g) \rbrace$ for all $f,g\in K[x]$.
    \end{enumerate}
\end{definition}

\begin{remark}
    Let $v$ be a pseudovaluation on $K[x]$.
    \begin{enumerate}[label=\textup{(\alph*)}]
        \item The definition implies that $\p_v \coloneqq v^{-1}(\infty) < K[x]$ is a prime ideal
            and that $v$ induces a valuation
            \begin{align}
                \bar{v}: K[x]/\p_v \to \hat{\RR}.
            \end{align}
            If $\p_v = (0)$, then $v=\bar{v}$ is a valuation in the usual sense.
        \item Let $\varphi\in K[x]$ be monic and irreducible, and let
            $L=K[x]/(\varphi)$.
            Then we have a one-to-one correspondence between the valuations $v_L: L \to \hat{\RR}$ extending $v_K$ and
            the pseudovaluations $v: K[x] \to \hat{\RR}$ extending $v_K$ with $v(\varphi)=\infty$.
            It is given by the mutually inverse maps
            \begin{align}
                v_L \mapsto \left( K[x] \twoheadrightarrow L \overset{v_L}{\to} \hat{\RR} \right)
            \end{align}
            and
            \begin{align}
                v \mapsto \left(L=K[x]/(\varphi) \overset{\bar{v}}{\longrightarrow} \hat{\RR}\right).
            \end{align}
    \end{enumerate}
\end{remark}

\begin{definition}
    A pseudovaluation $v: K[x] \to \widehat{\QQ}$ is called a \textit{Mac Lane pseudovaluation} if the following conditions hold.
    \begin{enumerate}[label=\textup{(\alph*)}]
        \setcounter{enumi}{2}
        \item $v|_K = v_K$,
        \item $v(x) \geq 0$,
        \item $v$ is discrete, i.e. the subgroup generated by the finite values,
            \begin{align}
                \Gamma_v\coloneqq\bigl\langle v(f):f\in K[x]\setminus\p_v\bigr\rangle,
            \end{align}
            is equal to $\frac{1}{e_v}\ZZ$ for some positive integer $e_v$,
        \item if $\p_v = (0)$, then the extension $k(v)|k$ has transcendence degree one.
    \end{enumerate}
    We denote the set of all Mac Lane pseudovaluations of $K[x]$ by $V(K[x])^*$ and the subset of Mac Lane valuations of $K[x]$ by $V(K[x]) \subset V(K[x])^*$.
\end{definition}

For the valuations in $V(K[x])$, there is a partial order: we say that $v \leq v'$ if and only if $v(f) \leq v'(f)$ for all $f\in K[x]$.

\begin{example}
    The Gauß valuation $v_{\scriptscriptstyle G} \in V(K[x])^*$ is minimal with respect to this partial order.
\end{example}

\begin{definition}
    Let $v\in V(K[x])$ be a Mac Lane valuation and let $f,g,h \in K[x]$.
    \begin{enumerate}[label=\textup{(\alph*)}]
        \item We say that $f$ and $g$ are \textit{$v$-equivalent} (written $f \sim_v g$) if $v(f-g) > v(f) = v(g)$.
        \item The polynomial $f$ is said to \textit{$v$-divide} $g$ (written $f\, |_v\, g$) if $g \sim_v fh$ for some $h \in K[x]$.
            It is called \textit{$v$-irreducible} if $f\, |_v\, gh$ implies $f\, |_v\, g$ or $f\, |_v\, h$,
            and \textit{$v$-minimal} if $f\, |_v\, g$ implies $\deg f \leq \deg g$.
        \item A polynomial $\phi \in K[x]$ is called a \textit{key polynomial} for $v$ if $\phi$ is monic, integral, $v$-irreducible and $v$-minimal.
    \end{enumerate}
\end{definition}

\begin{remark}
    One immediate consequence of these definitions is that any key polynomial $\phi$ for $v$ is irreducible (as an element of $K[x]$).
\end{remark}

\section{Augmentations and inductive valuations}\label{sec:augmentations-and-inductive-valuations}
Starting from a valuation and a key polynomial, augmentation produces a new pseudovaluation by prescribing a value for that polynomial which is at least its previous value. Iterating this construction leads to inductive valuations and their augmentation chains.

Let $\phi$ be a key polynomial for a valuation $v \in V(K[x])$ and let $\lambda \in \hat{\QQ}$ with $\lambda \geq v(\phi)$.
Then we can define a Mac Lane pseudovaluation
\begin{align} \label{eq:def-augmentation}
    v' = [v, v'(\phi) = \lambda] \in V(K[x])^*
\end{align}
as follows.
Any polynomial $f \in K[x]$ can be written uniquely as
\begin{align}
    f = \sum_i f_i \phi^i
\end{align}
with $\deg f_i < \deg \phi$. Then let
\begin{align}
    v'(f) \coloneqq \min_i(v(f_i) + i\cdot \lambda).
\end{align}

By construction, $v\leq v'$, and $v < v'$ if and only if $v(\phi) < \lambda$.

\begin{definition}
    The Mac Lane pseudovaluation $v'$ defined in~\eqref{eq:def-augmentation} is called the \textit{augmentation} of $v$ with respect to $(\phi,\lambda)$.
    If $v < v'$, $v'$ is called a \textit{proper augmentation}.
\end{definition}

\begin{definition}
    A Mac Lane pseudovaluation $v$ is called \textit{inductive} if there exists a chain of Mac Lane pseudovaluations
    \begin{align}
        v_{\scriptscriptstyle G} = v_0 \leq v_1 \leq \cdots \leq v_n = v,
    \end{align}
    such that $v_i$ is a proper augmentation of $v_{i-1}$ for $i=1,\dots, n$.
    We call $v_0,\dots,v_n$ the \textit{augmentation chain} representing $v$ and write \begin{align}
        v = [v_0,v_1(\phi_1) = \lambda_1,\dots, v_n(\phi_n) = \lambda_n].
    \end{align}
    Here $v_0 = v_{\scriptscriptstyle G}$, $v_n = v$, the polynomial $\phi_i$ is a key polynomial for $v_{i-1}$, $\lambda_i > v_{i-1}(\phi_i)$,
    and $v_i = [v_{i-1}, v_i(\phi_i) = \lambda_i]$.
    Also, $\lambda_i \neq \infty$ for $i < n$.
    We say that the augmentation chain $v_0,\dots,v_n$ is \textit{minimal} if
    \begin{align}
        \deg \phi_1 < \cdots < \deg \phi_n.
    \end{align}
\end{definition}


\section{Minimal augmentation chains}\label{sec:minimal-augmentation-chains}
The central structure result for these constructions describes Mac Lane pseudovaluations by minimal augmentation chains.

The following result combines the existence statements in~\cite[Theorem 4.31 and Corollary 4.67]{Rueth} with the uniqueness criterion in~\cite[Theorem 4.33]{Rueth}.

\begin{theorem}[Minimal augmentation chains]\label{thm:minimal-augmentation-chains}
    \begin{enumerate}[label=\textup{(\roman*)}]
        \item Every $v\in V(K[x])$ is inductive and can be represented by a minimal augmentation chain.
        \item If $K$ is complete, then every $v\in V(K[x])^*$ is inductive and can be represented by a minimal augmentation chain.
        \item Let
        \begin{align*}
            v_n&=[v_0,v_1(\phi_1)=\lambda_1,\dots,v_n(\phi_n)=\lambda_n],\\
            w_m&=[v_0,w_1(\psi_1)=\mu_1,\dots,w_m(\psi_m)=\mu_m]
        \end{align*}
        be inductive pseudovaluations represented by minimal augmentation chains.
        Then $v_n=w_m$ if and only if $m=n$ and, for every $i=1,\dots,n$,
        \begin{enumerate}[label=\textup{(\alph*)}]
            \item $\lambda_i=\mu_i$,
            \item $\deg\phi_i=\deg\psi_i$, and
            \item $v_{i-1}(\phi_i-\psi_i)\geq\lambda_i$.
        \end{enumerate}
    \end{enumerate}
\end{theorem}

\section{Discs, discoids, and valuations}\label{sec:discs-discoids-and-valuations}
We conclude by introducing discs and discoids in $\overline K$, which give a geometric description of subsets defined by valuation conditions on polynomials.

\begin{definition}
    \begin{enumerate}[label=\textup{(\alph*)}]
        \item Let $\alpha \in \overline K$ and let $\lambda \in \hat{\QQ}$.
            Then
                \begin{align}
                    D(\alpha, \lambda) \coloneqq \lbrace x \in \overline K:\, v_K(x - \alpha) \geq \lambda\rbrace
                \end{align}
            is the \textit{(closed) disc} with centre $\alpha$ and radius $\lambda$.
        \item Let $\phi \in K[x]$ be a monic irreducible polynomial and let $\lambda \in \hat{\QQ}$.
            The set
            \begin{align}
                D(\phi, \lambda) \coloneqq \lbrace x \in \overline K:\, v_K(\phi(x)) \geq \lambda \rbrace
            \end{align}
            is called the \textit{discoid} with centre $\phi$ and radius $\lambda$.
    \end{enumerate}
\end{definition}

\begin{remark}
    For the following description, assume that $K$ is complete, as in the cited source.
    Any disc with centre in $K$ is a discoid.
    Conversely, if $\phi = x - \alpha$ then $D(\phi, \lambda)$ coincides with the disc $D(\alpha, \lambda)$.
    More generally, let $\phi\in K[x]$ be monic and irreducible.
    Let $\alpha_1,\dots,\alpha_n\in\overline K$ be the roots of $\phi \in K[x]$, listed with multiplicity, so that $n=\deg\phi$,
    and let $\mu_1,\dots,\mu_n \in \hat{\QQ}$ be minimal\footnote{
        Our definition of discs implies that discs with large radius are smaller than those with small radius.
        For instance, the disc with centre $0\in K$ and radius $0$ is relatively large, as it consists of all integral elements.
    } such that $D(\alpha_i,\mu_i) \subset D(\phi,\lambda)$.
    Then
    \begin{align}
        D(\phi,\lambda) = \bigcup_{i=1}^n D(\alpha_i,\mu_i),
    \end{align}
    where the discs in this union need not be pairwise distinct; see \cite[Lemma 4.43]{Rueth}.
\end{remark}
